\chapter{方案或测试设计}

\section{总体思路}
这里描述总体设计思路、测试思路或实现路径。本章同时给出常用 \LaTeX{} 功能示例，便于观察图、表、公式、列表、脚注和交叉引用在当前模板中的排版效果。

\section{图表示例}

\subsection{流程图示例}
图~\ref{fig:test-flow-example} 展示了一个用 TikZ 绘制的测试流程图。该方式不依赖外部图片文件，适合绘制流程、结构或状态转换示意图。

\begin{figure}[H]
  \centering
  \begin{tikzpicture}[
    node distance=1.5cm,
    process/.style={rectangle, draw, rounded corners=2pt, minimum width=3cm, minimum height=0.9cm, align=center},
    decision/.style={diamond, draw, aspect=2.2, minimum width=3cm, align=center},
    arrow/.style={-{Latex[length=2.5mm]}, thick}
  ]
    \node[process] (prepare) {准备测试环境};
    \node[process, below=of prepare] (run) {执行测试用例};
    \node[decision, below=of run] (judge) {结果是否通过};
    \node[process, below left=1.3cm and 1.5cm of judge] (record) {记录测试结果};
    \node[process, below right=1.3cm and 1.5cm of judge] (fix) {定位并修正问题};

    \draw[arrow] (prepare) -- (run);
    \draw[arrow] (run) -- (judge);
    \draw[arrow] (judge) -- node[left] {是} (record);
    \draw[arrow] (judge) -- node[right] {否} (fix);
    \draw[arrow] (fix) |- (run);
  \end{tikzpicture}
  \caption{测试流程图示例}
  \label{fig:test-flow-example}
\end{figure}

\subsection{子图示例}
图~\ref{fig:subfig-example} 展示了子图排版效果，可用于对比不同状态、不同测试阶段或不同数据结果。

\begin{figure}[H]
  \centering
  \begin{subfigure}{0.44\textwidth}
    \centering
    \begin{tikzpicture}
      \draw[fill=gray!10] (0,0) rectangle (4,2.2);
      \draw[thick, -{Latex[length=2.2mm]}] (0.4,0.4) -- (3.6,0.4);
      \draw[thick, -{Latex[length=2.2mm]}] (0.4,0.4) -- (0.4,1.9);
      \draw[blue, thick] (0.4,0.6) -- (1.2,0.9) -- (2.0,1.1) -- (2.8,1.55) -- (3.5,1.75);
      \node at (2,2.0) {趋势 A};
    \end{tikzpicture}
    \caption{上升趋势}
  \end{subfigure}
  \hspace{0.04\textwidth}
  \begin{subfigure}{0.44\textwidth}
    \centering
    \begin{tikzpicture}
      \draw[fill=gray!10] (0,0) rectangle (4,2.2);
      \draw[thick, -{Latex[length=2.2mm]}] (0.4,0.4) -- (3.6,0.4);
      \draw[thick, -{Latex[length=2.2mm]}] (0.4,0.4) -- (0.4,1.9);
      \draw[red, thick] (0.4,1.6) -- (1.2,1.35) -- (2.0,1.0) -- (2.8,0.85) -- (3.5,0.55);
      \node at (2,2.0) {趋势 B};
    \end{tikzpicture}
    \caption{下降趋势}
  \end{subfigure}
  \caption{子图排版示例}
  \label{fig:subfig-example}
\end{figure}

\subsection{表格示例}
表~\ref{tab:test-case-example} 展示了三线表、多列合并和多行合并的效果，适合测试报告中的用例统计或结果汇总。

\begin{table}[H]
  \centering
  \caption{测试用例统计表示例}
  \label{tab:test-case-example}
  \begin{tabular}{llccc}
    \toprule
    测试阶段 & 测试类别 & 用例数量 & 通过数量 & 通过率 \\
    \midrule
    \multirow{2}{*}{功能验证}
      & 基础读写 & 12 & 12 & 100\% \\
      & 目录操作 & 8 & 8 & 100\% \\
    \midrule
    \multirow{2}{*}{可靠性验证}
      & 异常断电 & 10 & 10 & 100\% \\
      & 长时间运行 & 6 & 6 & 100\% \\
    \midrule
    \multicolumn{2}{r}{合计} & 36 & 36 & 100\% \\
    \bottomrule
  \end{tabular}
\end{table}

\section{其他常用功能示例}

\subsection{交叉引用与脚注}
模板支持图、表和公式的交叉引用。例如，流程图见图~\ref{fig:test-flow-example}，统计表见表~\ref{tab:test-case-example}，公式见式~\eqref{eq:pass-rate-example}。这里给出一个脚注示例\footnote{脚注可用于补充说明测试条件、数据来源或例外情况。}。

\subsection{公式示例}
通过率可按式~\eqref{eq:pass-rate-example} 计算：
\begin{equation}
  P = \frac{N_{\mathrm{pass}}}{N_{\mathrm{total}}} \times 100\%
  \label{eq:pass-rate-example}
\end{equation}

多行公式可使用 \texttt{align} 环境：
\begin{align}
  T_{\mathrm{avg}} &= \frac{1}{n}\sum_{i=1}^{n} T_i, \\
  T_{\mathrm{max}} &= \max\{T_1,T_2,\ldots,T_n\}.
\end{align}

\subsection{列表示例}
无序列表适合列出测试关注点：
\begin{itemize}
  \item 功能正确性；
  \item 性能稳定性；
  \item 异常工况恢复能力。
\end{itemize}

有序列表适合描述执行步骤：
\begin{enumerate}
  \item 初始化测试环境；
  \item 执行测试用例；
  \item 记录测试数据并形成结论。
\end{enumerate}

\subsection{引用与强调}
\begin{quote}
测试报告中的结论应由可复现的测试过程和可追溯的数据支撑。
\end{quote}

正文中可以使用 \textbf{加粗}、\emph{强调}、\texttt{等宽文本} 和 \colorbox{gray!15}{浅色底纹} 来突出不同类型的信息。

